%% ============================================================
%%  Reconocimiento de la Ejecución Biomecánica de Ejercicios
%%  mediante una Arquitectura Híbrida CNN-LSTM
%%  IEEE Conference Format — IEEEtran
%% ============================================================
\documentclass[conference]{IEEEtran}

\usepackage[utf8]{inputenc}
\usepackage[T1]{fontenc}
\usepackage[spanish]{babel}
\usepackage{amsmath}
\usepackage{amssymb}
\usepackage{graphicx}
\usepackage{booktabs}
\usepackage{cite}
\usepackage{hyperref}
\usepackage{url}
\usepackage{float}
\usepackage{array}

% ── Metadatos del documento ───────────────────────────────────
\hypersetup{
  pdftitle={Reconocimiento de la Ejecución Biomecánica de Ejercicios
            mediante una Arquitectura Híbrida CNN-LSTM},
  pdfauthor={Christopher Zúñiga, Adrian Arrieta Orozco},
  pdfkeywords={CNN, LSTM, biomecánica, reconocimiento de ejercicios,
               MediaPipe, aprendizaje profundo},
  colorlinks=true,
  linkcolor=black,
  citecolor=black,
  urlcolor=blue
}

% =============================================================
\begin{document}

% ── Título y autores ─────────────────────────────────────────
\title{Reconocimiento Computacional de la Ejecución Biomecánica
       de Ejercicios Mediante una Arquitectura Híbrida CNN-LSTM}

\author{
  \IEEEauthorblockN{Christopher Zúñiga}
  \IEEEauthorblockA{
    Escuela de Ciencias de la Computación e Informática\\
    Universidad de Costa Rica\\
    San José, Costa Rica\\
    \texttt{c.zuniga@ucr.ac.cr}
  }
  \and
  \IEEEauthorblockN{Adrian Arrieta Orozco}
  \IEEEauthorblockA{
    Escuela de Ciencias de la Computación e Informática\\
    Universidad de Costa Rica\\
    San José, Costa Rica\\
    \texttt{a.arrieta@ucr.ac.cr}
  }
}

\maketitle

% ── Abstract ─────────────────────────────────────────────────
\begin{abstract}
El análisis computacional del movimiento humano puede modelarse
como un problema espacio-temporal: la dimensión espacial codifica
la postura corporal en cada instante, mientras que la dimensión
temporal captura la dinámica del movimiento a lo largo del tiempo.
Este trabajo presenta el diseño de un sistema de reconocimiento
automático de la calidad de ejecución de ejercicios de fuerza,
clasificando la biomecánica observada como correcta o incorrecta.
Para ello, se seleccionaron tres ejercicios —sentadilla
(\textit{squat}), dominada (\textit{pull-up}) y peso muerto
(\textit{deadlift})— del conjunto de datos
\textit{Workout/Exercises Video Dataset} de Kaggle.
La arquitectura propuesta combina una red neuronal convolucional
(CNN) para la extracción de características espaciales por
fotograma y una red de memoria de corto y largo plazo (LSTM) para
modelar la evolución temporal del movimiento a lo largo de
secuencias de 30 fotogramas. Como representación de entrada, se
emplean los 33 puntos clave del esqueleto humano detectados
mediante el modelo BlazePose de MediaPipe, normalizados
estructuralmente y enriquecidos con ángulos articulares
biomecánicos calculados mediante álgebra vectorial. El enfoque
propuesto busca reducir la dependencia de procesamiento de imagen
cruda, enfocándose en variables cinemáticas invariantes que
faciliten la convergencia del modelo y su generalización en
condiciones de grabación variables.
\end{abstract}

\begin{IEEEkeywords}
aprendizaje profundo, biomecánica computacional, CNN-LSTM,
MediaPipe BlazePose, reconocimiento de actividad humana,
reconocimiento de ejercicios
\end{IEEEkeywords}

% =============================================================
\section{Introducción}
\label{sec:intro}

El análisis del movimiento humano es una línea de investigación
activa en visión por computadora y aprendizaje automático, con
aplicaciones en fisioterapia de rehabilitación, entrenamiento
deportivo de alto rendimiento, ergonomía laboral y sistemas de
asistencia inteligente~\cite{ren2020survey}. Dentro de este
dominio, el reconocimiento de la calidad de la ejecución de
ejercicios físicos representa un problema con implicaciones
directas en la prevención de lesiones biomecánicas.

El análisis del movimiento puede abordarse como un problema de
modelado espacio-temporal~\cite{ren2020survey}:

\begin{itemize}
  \item \textbf{Dimensión espacial:} posición y orientación de
        las articulaciones del cuerpo en cada instante.
  \item \textbf{Dimensión temporal:} dinámica y secuencia del
        movimiento a lo largo del tiempo.
\end{itemize}

Los enfoques basados en video ofrecen una representación natural
para capturar ambas dimensiones simultáneamente. Sin embargo, el
procesamiento directo de píxeles demanda recursos computacionales
elevados y puede ser sensible a variaciones de iluminación,
ángulo de cámara y características físicas del sujeto.

Una alternativa es estimar la pose del sujeto como un esqueleto
de puntos clave (\textit{landmarks}), abstrayendo la información
relevante antes de aplicar el modelo de clasificación. Herramientas
como MediaPipe BlazePose~\cite{mediapipe2024} permiten extraer en
tiempo real las coordenadas tridimensionales de 33 articulaciones
del cuerpo humano, proporcionando una representación compacta,
invariante a cambios de apariencia y apropiada para redes
neuronales secuenciales.

Este trabajo propone una arquitectura híbrida CNN-LSTM entrenada
sobre secuencias de landmarks y ángulos articulares para clasificar
la biomecánica de tres ejercicios de fuerza: \textit{squat},
\textit{pull-up} y \textit{deadlift}. La clasificación binaria
(ejecución correcta vs.\ incorrecta) permite un enfoque simplificado
que facilita la obtención de datos etiquetados y la evaluación
objetiva del modelo.

El resto del artículo se organiza como sigue: la
Sección~\ref{sec:related} presenta los trabajos relacionados;
la Sección~\ref{sec:dataset} describe el conjunto de datos y la
formulación del problema; la Sección~\ref{sec:methodology} detalla
el pipeline de procesamiento; la Sección~\ref{sec:architecture}
especifica la arquitectura propuesta; la
Sección~\ref{sec:design} describe el diseño orientado a objetos
del sistema; y la Sección~\ref{sec:conclusion} concluye el trabajo
y discute las líneas de trabajo futuro.

% =============================================================
\section{Trabajos Relacionados}
\label{sec:related}

\subsection{Redes Neuronales Convolucionales para Datos Visuales}

Las redes neuronales convolucionales (CNN) se distinguen del resto
de arquitecturas de aprendizaje profundo por su distribución
tridimensional de neuronas, lo que las vuelve eficientes para el
procesamiento de datos visuales. Su estructura puede describirse
mediante tres componentes fundamentales:

\begin{enumerate}
  \item \textbf{Capa convolucional:} procesa regiones localizadas
        de la entrada mediante filtros que extraen características
        como bordes, texturas y patrones estructurales.
  \item \textbf{Capa de agrupación (\textit{pooling}):} reduce las
        dimensiones espaciales preservando la información esencial,
        reduciendo la complejidad computacional.
  \item \textbf{Capa completamente conectada:} clasifica las
        características extraídas mediante propagación hacia
        adelante y retropropagación.
\end{enumerate}

\subsection{Redes LSTM para Modelado Temporal}

Las redes de memoria de corto y largo plazo (LSTM) son un tipo de
red neuronal recurrente (RNN) diseñadas para mantener dependencias
a largo plazo en secuencias. Cada unidad LSTM incorpora una celda
de memoria controlada por tres compuertas~\cite{ullah2018action}:

\begin{itemize}
  \item \textbf{Compuerta de entrada (\textit{input gate}):}
        decide cuánta información nueva se almacena en memoria.
  \item \textbf{Compuerta de salida (\textit{output gate}):}
        regula cuánta información se propaga a la capa siguiente.
  \item \textbf{Compuerta de olvido (\textit{forget gate}):}
        controla qué información previa se descarta de la memoria.
\end{itemize}

Este mecanismo permite a la LSTM modelar dependencias temporales
de largo alcance, como la evolución progresiva de un movimiento
a lo largo de múltiples fotogramas.

\subsection{Arquitecturas Híbridas CNN-LSTM para Reconocimiento
            de Actividad}

Ullah et al.~\cite{ullah2018action} demostraron que la combinación
de características espaciales extraídas por una CNN con el modelado
temporal de una LSTM mejora significativamente la precisión en el
reconocimiento de acciones en secuencias de video. En esta
arquitectura, la CNN actúa como extractor de características por
fotograma, y la LSTM modela la secuencia de vectores resultantes
para capturar la dinámica del movimiento.

Esta combinación resulta especialmente apropiada para el análisis
biomecánico, donde la postura en un instante aislado es insuficiente
para evaluar la calidad del movimiento: se requiere analizar cómo
las articulaciones se desplazan y coordinan a lo largo del tiempo.

\subsection{Representación Esqueletal y Análisis Biomecánico}

Ren et al.~\cite{ren2020survey} validan protocolos de
preprocesamiento esenciales para el reconocimiento de acción
basado en esqueleto, incluyendo la normalización por traslación
centrada en la cadera (\textit{hip-center translation}) y el
escalado por longitud de torso para lograr invariancia espacial.
Winter~\cite{winter2009biomechanics} establece que el análisis del
movimiento humano debe fundamentarse en vectores de fuerza y
ángulos articulares, lo que justifica la ingeniería de
características biomecánicas como etapa de preprocesamiento.
Zhou et al.~\cite{zhou2019continuity} respaldan el uso de
representaciones basadas en ángulos y rotaciones para mejorar la
eficiencia del aprendizaje en arquitecturas de aprendizaje profundo.

% =============================================================
\section{Conjunto de Datos y Formulación del Problema}
\label{sec:dataset}

\subsection{Conjunto de Datos}

El conjunto de datos seleccionado es el \textit{Workout/Exercises
Video Dataset}~\cite{kaggle_dataset}, disponible en Kaggle. Este
dataset contiene aproximadamente 652 videos distribuidos en 25
categorías de ejercicio. Para el presente trabajo se seleccionaron
tres categorías:

\begin{table}[htbp]
  \caption{Distribución del subconjunto de datos utilizado}
  \label{tab:dataset}
  \centering
  \begin{tabular}{lcc}
    \toprule
    \textbf{Ejercicio} & \textbf{Videos disponibles} & \textbf{Formato} \\
    \midrule
    Deadlift  & 32 & .mp4 \\
    Squat     & 29 & .mp4 / .MOV \\
    Pull-up   & 26 & .mp4 \\
    \midrule
    \textbf{Total} & \textbf{87} & — \\
    \bottomrule
  \end{tabular}
\end{table}

\subsection{Formulación del Problema}

El problema se formula como una clasificación binaria. Dado un
video $V$ que contiene la ejecución de uno de los tres ejercicios,
el objetivo del modelo es determinar si la biomecánica del
movimiento observado corresponde a una ejecución correcta o
incorrecta:

\begin{equation}
  \hat{y} = f(X_{1:T}), \quad \hat{y} \in \{0, 1\}
  \label{eq:problem}
\end{equation}

donde $X_{1:T}$ representa la secuencia de características
extraídas sobre $T = 30$ fotogramas del video, y $\hat{y} = 1$
corresponde a ejecución correcta.

La reducción a tres ejercicios específicos y la formulación binaria
del problema persiguen aumentar la precisión del modelo al
enfocarse en los patrones biomecánicos particulares de cada
movimiento, en lugar de abordar un problema de clasificación
multi-ejercicio de mayor complejidad.

% =============================================================
\section{Metodología de Procesamiento de Datos}
\label{sec:methodology}

\subsection{Extracción de Fotogramas y Detección de Landmarks}

Para cada video del conjunto de datos, se extrae una secuencia
fija de $T = 30$ fotogramas mediante muestreo uniforme a lo largo
de la duración del video. Sobre cada fotograma se aplica el modelo
BlazePose de MediaPipe~\cite{mediapipe2024}, el cual detecta 33
puntos clave (\textit{landmarks}) del cuerpo humano con sus
coordenadas tridimensionales $(x, y, z)$.

\subsection{Normalización Estructural}

Para que el modelo sea invariante a la distancia de la cámara y
a la constitución física del sujeto, se aplica una normalización
estructural en dos pasos~\cite{ren2020survey}:

\begin{enumerate}
  \item \textbf{Traslación:} los landmarks se traducen para que la
        cadera (\textit{hip-center}) quede en el origen del
        sistema de coordenadas.
  \item \textbf{Escalado:} las coordenadas se normalizan por la
        longitud del torso del sujeto, eliminando la dependencia
        de la escala corporal.
\end{enumerate}

\subsection{Ingeniería de Características Biomecánicas}

Las coordenadas de los landmarks se complementan con ángulos
articulares relevantes para cada ejercicio. El ángulo $\theta$
entre dos segmentos óseos representados como vectores $\mathbf{A}$
y $\mathbf{B}$ se calcula mediante el producto punto de sus
vectores unitarios~\cite{winter2009biomechanics, zhou2019continuity}:

\begin{equation}
  \theta = \arccos\!\left(\frac{\mathbf{A} \cdot \mathbf{B}}
           {\|\mathbf{A}\| \, \|\mathbf{B}\|}\right)
  \label{eq:angle}
\end{equation}

Complementariamente, se emplea la ley del coseno para derivar
ángulos internos en triángulos formados por tres articulaciones.
Los ángulos articulares calculados por ejercicio incluyen:

\begin{itemize}
  \item \textbf{Squat:} ángulo de rodilla y de cadera.
  \item \textbf{Pull-up:} ángulo de codo y de hombro.
  \item \textbf{Deadlift:} ángulo de cadera, rodilla y columna
        lumbar (estimada).
\end{itemize}

Este enfoque de ingeniería de características reduce la
dimensionalidad de los datos de entrada, elimina redundancias
espaciales y permite que la red neuronal converja sobre variables
biomecánicas invariantes~\cite{zhou2019continuity}.

\subsection{Representación como Tensores}

Los datos procesados se organizan en tensores de la
forma~\cite{w3schools_tensor}:

\begin{equation}
  \mathbf{X} \in \mathbb{R}^{B \times T \times F}
  \label{eq:tensor}
\end{equation}

donde $B$ es el tamaño del lote (\textit{batch}), $T = 30$ es el
número de fotogramas y $F$ es la dimensión de características por
fotograma, compuesta por las coordenadas $(x, y, z)$ de los 33
landmarks más los ángulos biomecánicos calculados.

Para evitar sesgos en el entrenamiento, se aplica balanceo de
clases garantizando que cada lote contenga la misma proporción
de ejecuciones correctas e incorrectas.

% =============================================================
\section{Arquitectura del Modelo}
\label{sec:architecture}

Se propone una arquitectura híbrida CNN-LSTM que procesa la
dinámica espacio-temporal del movimiento en dos etapas, como
se ilustra a continuación.

\subsection{Extracción Espacial mediante CNN}

Se aplican capas convolucionales de una dimensión (Conv1D) sobre
el vector de características de cada fotograma. Su función es
detectar patrones locales de postura, como la alineación de la
espalda o la profundidad del movimiento. Cada capa convolucional
extrae un vector de características comprimido por fotograma.

\subsection{Análisis Temporal mediante LSTM}

Los vectores de características obtenidos por la CNN en cada
fotograma se organizan en una secuencia temporal y se inyectan
en una red LSTM. Las celdas de memoria y compuertas de la LSTM
identifican la evolución del ejercicio a lo largo de los 30
fotogramas, detectando anomalías en la velocidad, profundidad
o ritmo de ejecución.

\subsection{Clasificación Final}

Una capa densa completamente conectada con activación sigmoide
genera una probabilidad $p \in [0, 1]$ de ejecución correcta,
donde:

\begin{equation}
  \hat{y} =
  \begin{cases}
    1 & \text{si } p \geq 0.5 \quad \text{(ejecución correcta)} \\
    0 & \text{si } p < 0.5  \quad \text{(ejecución incorrecta)}
  \end{cases}
  \label{eq:decision}
\end{equation}

% =============================================================
\section{Diseño del Sistema}
\label{sec:design}

Para gestionar la complejidad del pipeline y mantener separación
de responsabilidades, el sistema se implementa mediante una
arquitectura orientada a objetos con tres módulos principales:

\begin{itemize}
  \item \textbf{DataManager:} transforma los videos crudos en
        flujos de datos estructurados. Encapsula la extracción
        de landmarks con MediaPipe, la normalización estructural
        y el cálculo de ángulos articulares.

  \item \textbf{DataLoader:} actúa como puente entre los datos
        procesados y la red neuronal. Organiza los datos en
        tensores, aplica el balanceo de clases y entrega lotes
        normalizados y equitativos durante el entrenamiento.

  \item \textbf{RecognitionModel:} define la arquitectura
        CNN-LSTM, gestiona el entrenamiento, el ajuste de
        hiperparámetros y la evaluación del rendimiento.
\end{itemize}

Esta separación permite modificar cualquier etapa del pipeline
de forma independiente, facilitando la experimentación con
diferentes configuraciones de arquitectura o estrategias de
preprocesamiento.

% =============================================================
\section{Conclusión y Trabajo Futuro}
\label{sec:conclusion}

Este trabajo presenta el diseño de un sistema de reconocimiento
automático de la calidad de ejecución biomecánica de ejercicios
de fuerza, fundamentado en una arquitectura híbrida CNN-LSTM y
en la extracción de representaciones esqueletales mediante
MediaPipe BlazePose. La formulación del problema como
clasificación binaria, combinada con la normalización estructural
y la ingeniería de características biomecánicas, proporciona
una base sólida para el entrenamiento de un modelo generalizable
a condiciones de grabación variables.

Como trabajo futuro, se contempla la recolección y etiquetado
de datos de ejecución correcta e incorrecta, la implementación
completa del pipeline de preprocesamiento, el entrenamiento y
ajuste de hiperparámetros del modelo, y la evaluación de su
rendimiento mediante métricas estándar de clasificación
(\textit{accuracy}, precisión, recall, F1-score) y matrices
de confusión por clase.

% =============================================================
\section*{Reconocimientos}

Los autores agradecen al personal docente del curso de
Inteligencia Artificial, Ciclo IV, de la Universidad de Costa
Rica por la orientación brindada durante el desarrollo de este
proyecto.

% =============================================================
\bibliographystyle{IEEEtran}
\bibliography{references}

\end{document}
